\section*{الفصل \ref{ch:g7:literal} --- العبارات الحرفية}

\begin{solution}{exo:g7:literal:1}
$7x$; \quad $3y$; \quad $5ab$; \quad $x^2$; \quad $2(a + 4)$.
\end{solution}

\begin{solution}{exo:g7:literal:2}
$5ab = 5 \times a \times b$; \quad
$3(x+2) = 3 \times (x + 2)$; \quad
$x^2 + 4x = x \times x + 4 \times x$; \quad
$2a^3 = 2 \times a \times a \times a$.
\end{solution}

\begin{solution}{exo:g7:literal:3}
من أجل $x = 3$: \ $5 \times 3 + 1 = 16$؛ \quad و $3 \times 3 = 9$؛ \quad
و $2 \times 9 - 3 = 15$؛ \quad و $10 - 6 = 4$.
\end{solution}

\begin{solution}{exo:g7:literal:4}
$3a + 2b = 12 + 3 = 15$; \qquad
$ab + 5 = 4 \times 1.5 + 5 = 6 + 5 = 11$.
\end{solution}

\begin{solution}{exo:g7:literal:5}
$8x + 5x = 13x$; \quad
$9a - 4a + a = 6a$; \quad
$6y + 2 + 3y + 7 = 9y + 9$; \quad
$5t + 3 - 2t - 3 = 3t$.
\end{solution}

\begin{solution}{exo:g7:literal:6}
المربع: $P = 4c$. \omterm{def:g6:shapes:triangles}{والمثلث المتساوي الساقين}: $P = b + 2s$. والأرغفة:
$p = 1.20\,n$ يورو.
\end{solution}

\begin{solution}{exo:g7:literal:7}
$2n$: الضعف؛ \ و $n + 2$: أكثر من $n$ باثنين؛ \ و $n^2$: المربع؛
\ و $\frac n2$: النصف.
\end{solution}

\begin{solution}{exo:g7:literal:8}
$3(x+4) = 3x + 12$: \emph{صواب} من أجل كل $x$ (بالتوزيعية ---
وهذا برهان، لا حاجة إلى اِختبار).

$2x + 5 = 7x$: اِختبر $x = 2$: اليسار $9$، واليمين $14$ --- \emph{خطأ}
عمومًا (فهو لا يصحّ إلا من أجل $x = 1$).

$x + x = x^2$: اِختبر $x = 3$: اليسار $6$، واليمين $9$ --- \emph{خطأ}
عمومًا (إذ $x + x = 2x$).
\end{solution}

\begin{solution}{exo:g7:literal:9}
\emph{1.} الطول: $w + 3$. \omterm{def:g6:measure:perimeter}{والمحيط}:
\[
P = 2\bigl(w + (w + 3)\bigr) = 2(2w + 3) = 4w + 6 .
\]

\emph{2.} من أجل $w = 5$: مباشرةً، الضلعان $5$ و $8$، إذن
$P = 2 \times 13 = 26$ سم؛ وبالصيغة، $4 \times 5 + 6 = 26$ سم.
وهما متوافقان.
\end{solution}

\begin{solution}{exo:g7:literal:10}
باتّباع الخطوات بالحرف $n$:
\[
n \ \to\ n + 6 \ \to\ 3(n + 6) = 3n + 18 \ \to\ 3n + 18 - 18 = 3n
\ \to\ \frac{3n}{n} = 3 .
\]
فيحصل الجميع على $3$.
\end{solution}

\begin{solution}{exo:g7:literal:11}
على صورة \omterm{ex:g1:subtraction:difference}{فرق}: $A = a^2 - b^2$. وبقطع حرف L إلى مستطيلين:
شريط بعرض كامل اِرتفاعه $a - b$، وفوقه
مستطيل أضيق:
\[
A = a(a - b) + b(a - b)
\]
(المستطيل السفلي عرضه $a$ واِرتفاعه $a - b$، والمستطيل العلوي عرضه $a - b$
واِرتفاعه $b$ يعطي الصورة $a(a-b) + (a-b)b$ --- وهي نفسها).
ومن أجل $a = 5$ و $b = 2$: $a^2 - b^2 = 25 - 4 = 21$، و
$5 \times 3 + 2 \times 3 = 15 + 6 = 21$. فالعبارتان متوافقتان ---
وهما كذلك دائمًا: فهذه هي المتطابقة $a^2 - b^2 = (a + b)(a - b)$ في
\cref{ch:g9:algebra} متنكّرةً.
\end{solution}

\begin{solution}{pb:g7:literal:1}
\textbf{1.} الأشكال من $1$ إلى $4$ تحتاج $4$ و $7$ و $10$ و $13$
عودًا.

\textbf{2.} المربع الجديد يشترك في ضلع واحد مع الشكل
السابق، فلا يكون جديدًا إلا $3$ من أضلاعه $4$. واِنطلاقًا من
الشكل $4$ ($13$ عودًا)، تكلّف ستة مربعات أخرى
$6 \times 3 = 18$: فالشكل $10$ يحتاج $13 + 18 = 31$ عودًا.

\textbf{3.} بقراءة الصفّ من اليسار: يبدأ بعود
عمودي واحد، ثم يسهم كل واحد من المربعات $n$ بضلع
أعلى وضلع أسفل وضلع أيمن --- أي $3$ أعواد لكل مربع. \omterm{def:g1:addition:def}{والمجموع}:
$3n + 1$؛ و«$+1$» هي أول عود عمودي. والتحقّق:
$n = 1, 2, 3, 4$ تعطي $4, 7, 10, 13$.

\textbf{4.} $4 + 3(n - 1) = 4 + 3n - 3 = 3n + 1$: فصيغة أمير
تنتشر إلى العبارة المختزلة نفسها بالضبط. فالتوافق
تامّ.

\textbf{5.} لينا: $2n + (n + 1) = 3n + 1$ --- وهي نفسها من جديد.
ثلاث طرق للرؤية، وصورة مختزلة واحدة. والشكل $100$:
$3 \times 100 + 1 = 301$ عودًا.

\textbf{6.} الأعداد: $3$ و $5$ و $7$ أعواد. وكل مثلث جديد
يتّكئ على ضلع موجود ولا يضيف إلا $2$ عودين، اِنطلاقًا
من $3$: فالشكل $n$ يحتاج $3 + 2(n - 1) = 2n + 1$ عودًا.

\textbf{7.} اِعكس $2n + 1 = 49$: فنزع العود الأول يترك
$48$، ثم يقابل كل مثلث $2$:
$48 \div 2 = 24$ مثلثًا.

\textbf{8.} في صفّ من $n$ طاولة، تقدّم كل طاولة ضلعها
الأعلى والأسفل ($2n$ مقعدًا)، وتقدّم الطاولتان الطرفيتان ضلعًا
إضافيًّا لكل منهما ($+2$): أي $2n + 2$ مقعدًا. والتحقّق: $4$ و $6$ و $8$
من أجل $n = 1, 2, 3$.

\textbf{9.} صفّ واحد: $2n + 2 \geq 30$ يحتاج $2n \geq 28$، إذن
$14$ طاولة. وبقسمة الطاولات $14$ نفسها إلى صفّين من
$7$: يجلس في كل صفّ $2 \times 7 + 2 = 16$، ومعًا $32$ ---
أي اِثنان أكثر من الصفّ الواحد، لأن كل صفّ جديد يجلب طرفين
جديدين. (وكل قسمة تكسب $2$ مقعدين؛ والمُطعِمون يعرفون ذلك.)

\textbf{10.} حرف L ذو \omterm{def:g6:measure:volume}{الحجم} $n$ فيه $n$ نقطة صاعدة و $n$
نقطة عرضية، مع اِشتراك نقطة الركن: أي $n + n - 1 = 2n - 1$ نقطة.
والأشكال $1$ و $2$ و $3$: $1$ و $3$ و $5$ نقط --- وهي \emph{\omterm{def:g3:numbers:evenodd}{الأعداد
الفردية}}. وبتعشيش كل واحد حول الذي يليه، تملأ الحروف
مربع نقط $n \times n$: وهي قصة
\cref{pb:g8:equations:1}.

\textbf{11.} $3(n + 2) - 5 = 3n + 6 - 5 = 3n + 1$: فالادّعاء
يُختزل إلى الصيغة المبرهن عليها، إذن هو صحيح من أجل كل $n$ ---
ثبت بالجبر، لا بالأمثلة.

\textbf{12.} الاختبارات: $4 \geq 2$ و $9 \geq 3$ و $100 \geq 10$: وكلها
تنجح. لكن $n = 0.5$ تعطي $n^2 = 0.25 < 0.5$: فيفشل الادّعاء.
ولم تبرهن الاختبارات الثلاثة على شيء بخصوص \emph{كل} الأعداد (فهي
لم تفشل إلا في إيجاد مشكلة)؛ أما المثال المضادّ الواحد
\emph{فيبرهن} على أن الادّعاء العامّ خطأ
(\cref{met:g7:literal:test}). فالحرف يمثّل كل عدد
--- \omterm{def:g6:decimals:places}{والأعداد العشرية} منها.

\textbf{13.} المناطق: $2$ نقطتان $\to 2$؛ و $3$ نقط $\to 4$؛
و $4$ نقط $\to 8$؛ و $5$ نقط $\to 16$. ويهمس النمط
بأن «كل نقطة تضاعف العدد»: فتُتوقَّع $32$ منطقة من أجل
$6$ نقط.

\textbf{14.} العدّ الدقيق يعطي $31$ منطقة --- لا $32$.
فتخمين المضاعفة ميّت: نجا من أربعة فحوص
وفشل في الخامس. (وهذا المثال المضادّ مشهور بما يكفي ليكون له
اِسم، «فخّ مناطق الدائرة»؛ ولا اِختصارات بتباعد منتظم
--- فتلاقي ثلاثة أوتار في نقطة واحدة يدمج المناطق ويفسد العدّ.)

\textbf{15.} كان البرهان يحتاج سببًا \emph{بنيويًّا}
لكون إضافة نقطة تضاعف المناطق --- ولا سبب
كهذا: فالمضاعفة كانت مصادفة في الحالات الصغيرة.
وأما صيغ أعواد الثقاب فكانت آمنة لأن كل واحدة قُرئت من
البناء نفسه (فكل مربع يكلّف ظاهرًا $3$ أعواد،
وكل طاولة تقدّم ظاهرًا $2$ مقعدين): بنية لا اِختبارًا.
والتوقّع المبرهن عليه: $3 \times 2\,026 + 1 = 6\,079$ عودًا.
\end{solution}
